% =====================================================================
%  ANNEXE B — 600 PHRASES PRÊTES À L'EMPLOI
% =====================================================================
\chapter{600 phrases prêtes à l'emploi}

Voici votre \textbf{boîte à outils}. Ces phrases ne s'étudient pas comme du vocabulaire :
elles se \textbf{mémorisent comme des blocs}, entiers, avec leur son. Une phrase apprise
par c\oe ur est disponible instantanément à l'oral --- alors qu'une phrase construite mot à
mot coûte trois secondes de réflexion.

\begin{regle}[Comment utiliser ce répertoire]
\textbf{1.} Choisissez \textbf{dix phrases par jour}, pas plus.\\
\textbf{2.} Lisez-les à voix haute trois fois, en exagérant l'intonation.\\
\textbf{3.} Enregistrez-vous, écoutez, corrigez.\\
\textbf{4.} Utilisez chacune d'elles \textbf{au moins une fois dans la journée} --- même en
vous parlant à vous-même sous la douche.\\
\textbf{5.} Le week-end, test : cachez la colonne française, lisez la colonne anglaise.
Puis cachez l'anglais, dites la phrase en partant du français.\\
Une phrase est « acquise » quand vous pouvez la dire en moins de deux secondes, sans
traduire. Comptez environ \textbf{40 jours} pour maîtriser les 600.
\end{regle}

\section{1. Saluer et prendre congé}
\begin{tableau}[30 formules de salutation]
\footnotesize
\begin{longtable}{@{}p{8.0cm}p{6.6cm}@{}}
\entete{Anglais} & \entete{Français} \tabularnewline
\endfirsthead
\entete{Anglais} & \entete{Français} \tabularnewline
\endhead
\phr{Hello! How are you?}{Bonjour ! Comment allez-vous ?}
\phr{Hi there! Long time no see.}{Salut ! Ça fait longtemps.}
\phr{Good morning, everyone.}{Bonjour à tous.}
\phr{Good afternoon. How's it going?}{Bonjour. Comment ça va ?}
\phr{Good evening. Nice to see you again.}{Bonsoir. Content de vous revoir.}
\phr{Hey! What's up?}{Salut ! Quoi de neuf ?}
\phr{Hi! How have you been?}{Salut ! Comment vas-tu depuis la dernière fois ?}
\phr{Nice to meet you.}{Enchanté de faire votre connaissance.}
\phr{It's a pleasure to meet you.}{C'est un plaisir de vous rencontrer.}
\phr{I've heard so much about you.}{J'ai beaucoup entendu parler de vous.}
\phr{Sorry I'm late.}{Désolé, je suis en retard.}
\phr{Welcome! Come in.}{Bienvenue ! Entrez.}
\phr{How do you do?}{Comment allez-vous ? (très formel, première rencontre)}
\phr{Good to see you.}{Content de te voir.}
\phr{How's life treating you?}{Comment va la vie ?}
\phr{See you later!}{À plus tard !}
\phr{See you tomorrow morning.}{À demain matin.}
\phr{Take care!}{Prends soin de toi !}
\phr{Have a nice day!}{Bonne journée !}
\phr{Have a good weekend.}{Bon week-end.}
\phr{Have a safe trip.}{Bon voyage / bonne route.}
\phr{I'd better get going.}{Il faut que j'y aille.}
\phr{I must be off --- time flies.}{Je dois y aller --- le temps passe vite.}
\phr{It was lovely talking to you.}{C'était vraiment agréable de parler avec vous.}
\phr{Thanks for having me.}{Merci de m'avoir reçu.}
\phr{Give my regards to your family.}{Transmettez mes amitiés à votre famille.}
\phr{Say hello to Aïcha for me.}{Dis bonjour à Aïcha de ma part.}
\phr{Let's keep in touch.}{Restons en contact.}
\phr{Catch you later!}{On se capte plus tard ! (familier)}
\phr{Goodbye, and good luck with everything.}{Au revoir, et bonne chance pour tout.}
\end{longtable}
\end{tableau}

\section{2. Se présenter}
\begin{tableau}[28 phrases pour parler de soi]
\footnotesize
\begin{longtable}{@{}p{8.0cm}p{6.6cm}@{}}
\entete{Anglais} & \entete{Français} \tabularnewline
\endfirsthead
\entete{Anglais} & \entete{Français} \tabularnewline
\endhead
\phr{My name is Moussa --- but everyone calls me Mouss.}{Je m'appelle Moussa --- mais tout le monde m'appelle Mouss.}
\phr{I'm from Senegal, from a town called Thiès.}{Je viens du Sénégal, d'une ville qui s'appelle Thiès.}
\phr{I was born in 1995.}{Je suis né en 1995.}
\phr{I'm in my late twenties.}{J'ai la fin de la vingtaine.}
\phr{I live in Dakar with my family.}{J'habite à Dakar avec ma famille.}
\phr{I've been living here for five years.}{J'habite ici depuis cinq ans.}
\phr{I work as an accountant.}{Je travaille comme comptable.}
\phr{I've been working in this company since 2020.}{Je travaille dans cette entreprise depuis 2020.}
\phr{I'm a student at university.}{Je suis étudiant à l'université.}
\phr{I'm studying business administration.}{J'étudie la gestion.}
\phr{I'm currently looking for a job.}{Je cherche un emploi en ce moment.}
\phr{I'm self-employed.}{Je travaille à mon compte.}
\phr{In my free time, I play football and read.}{Pendant mon temps libre, je joue au football et je lis.}
\phr{I'm really into photography.}{Je suis passionné de photographie.}
\phr{I've been learning English for a year.}{J'apprends l'anglais depuis un an.}
\phr{I also speak French and Wolof.}{Je parle aussi français et wolof.}
\phr{I'm married and I have two children.}{Je suis marié et j'ai deux enfants.}
\phr{I'm single at the moment.}{Je suis célibataire en ce moment.}
\phr{I come from a big family --- we are seven.}{Je viens d'une grande famille --- nous sommes sept.}
\phr{My dream is to work abroad.}{Mon rêve est de travailler à l'étranger.}
\phr{I'd describe myself as calm and curious.}{Je me décrirais comme calme et curieux.}
\phr{People say I'm a good listener.}{On dit que je sais écouter.}
\phr{What about you? Tell me about yourself.}{Et vous ? Parlez-moi de vous.}
\phr{And what do you do for a living?}{Et que faites-vous dans la vie ?}
\phr{How long have you lived here?}{Depuis combien de temps habitez-vous ici ?}
\phr{Where are you from originally?}{D'où venez-vous à l'origine ?}
\phr{Sorry, how do you pronounce your name?}{Pardon, comment prononce-t-on votre nom ?}
\phr{Could you spell that for me, please?}{Pouvez-vous me l'épeler, s'il vous plaît ?}
\end{longtable}
\end{tableau}

\section{3. Répondre, réagir, montrer qu'on écoute}
\begin{tableau}[30 réactions --- indispensables pour ne pas paraître froid]
\footnotesize
\begin{longtable}{@{}p{8.0cm}p{6.6cm}@{}}
\entete{Anglais} & \entete{Français} \tabularnewline
\endfirsthead
\entete{Anglais} & \entete{Français} \tabularnewline
\endhead
\phr{I'm fine, thanks. And you?}{Je vais bien, merci. Et vous ?}
\phr{Not bad, thanks --- you?}{Pas mal, merci --- et toi ?}
\phr{Pretty good, actually.}{Plutôt bien, en fait.}
\phr{Can't complain.}{Je ne peux pas me plaindre.}
\phr{I'm exhausted --- it's been a long week.}{Je suis épuisé --- la semaine a été longue.}
\phr{Things are going well, thanks.}{Les choses se passent bien, merci.}
\phr{Really? Tell me more.}{Vraiment ? Racontez-moi.}
\phr{That's interesting.}{C'est intéressant.}
\phr{Wow, that's amazing!}{Waouh, c'est génial !}
\phr{That sounds great.}{Ça a l'air super.}
\phr{I'm so happy for you.}{Je suis vraiment content pour vous.}
\phr{Congratulations! Well done.}{Félicitations ! Bravo.}
\phr{That's terrible --- I'm so sorry.}{C'est terrible --- je suis désolé.}
\phr{Oh no, what happened?}{Oh non, qu'est-ce qui s'est passé ?}
\phr{That must have been difficult.}{Ça a dû être difficile.}
\phr{I understand how you feel.}{Je comprends ce que vous ressentez.}
\phr{Are you all right?}{Vous allez bien ?}
\phr{Take your time.}{Prenez votre temps.}
\phr{No worries.}{Pas de souci.}
\phr{It doesn't matter.}{Ça ne fait rien.}
\phr{Fair enough.}{C'est juste / je l'admets.}
\phr{I see what you mean.}{Je vois ce que vous voulez dire.}
\phr{That makes sense.}{C'est logique.}
\phr{Right, I get it.}{D'accord, je comprends.}
\phr{You're joking!}{Vous plaisantez !}
\phr{No way!}{Pas possible !}
\phr{That's unbelievable.}{C'est incroyable.}
\phr{I can imagine.}{J'imagine.}
\phr{Same here.}{Pareil pour moi.}
\phr{Me too --- exactly the same.}{Moi aussi --- exactement pareil.}
\end{longtable}
\end{tableau}

\section{4. Phrases de survie : quand vous ne comprenez pas}
\begin{tableau}[28 phrases --- les plus importantes de tout ce répertoire]
\footnotesize
\begin{longtable}{@{}p{8.0cm}p{6.6cm}@{}}
\entete{Anglais} & \entete{Français} \tabularnewline
\endfirsthead
\entete{Anglais} & \entete{Français} \tabularnewline
\endhead
\phr{Sorry, could you say that again, please?}{Pardon, pouvez-vous répéter, s'il vous plaît ?}
\phr{Could you speak more slowly, please?}{Pourriez-vous parler plus lentement ?}
\phr{I'm sorry, I didn't catch that.}{Désolé, je n'ai pas saisi.}
\phr{What does that mean?}{Qu'est-ce que ça veut dire ?}
\phr{How do you say ``vélo'' in English?}{Comment dit-on « vélo » en anglais ?}
\phr{What's the word for this?}{Comment appelle-t-on cela ?}
\phr{Could you spell it for me, please?}{Pouvez-vous me l'épeler, s'il vous plaît ?}
\phr{Sorry, could you write it down?}{Pardon, pourriez-vous l'écrire ?}
\phr{Do you mean\ldots{} ?}{Vous voulez dire\ldots{} ?}
\phr{So you're saying that\ldots{} ?}{Donc vous dites que\ldots{} ?}
\phr{Let me see if I understand correctly.}{Voyons si j'ai bien compris.}
\phr{I understand the words, but not the meaning.}{Je comprends les mots, mais pas le sens.}
\phr{My English isn't perfect yet, but I'm working on it.}{Mon anglais n'est pas encore parfait, mais je progresse.}
\phr{Could you use a simpler word, please?}{Pourriez-vous employer un mot plus simple ?}
\phr{I'm not sure I follow you.}{Je ne suis pas sûr de vous suivre.}
\phr{Sorry, I've lost the thread.}{Désolé, j'ai perdu le fil.}
\phr{Could you give me an example?}{Pourriez-vous me donner un exemple ?}
\phr{Is that a common expression?}{Est-ce une expression courante ?}
\phr{Would you mind explaining that again?}{Cela vous ennuierait-il d'expliquer encore ?}
\phr{I know the word but I can't pronounce it.}{Je connais le mot mais je n'arrive pas à le prononcer.}
\phr{Please stop me if I make a mistake.}{Arrêtez-moi si je fais une erreur.}
\phr{Can we speak a bit more slowly? I'm still learning.}{Pouvons-nous parler un peu plus lentement ? J'apprends encore.}
\phr{Just a second --- let me think.}{Un instant --- laissez-moi réfléchir.}
\phr{Give me a moment, please.}{Donnez-moi un instant, s'il vous plaît.}
\phr{What did you say the name was?}{Quel était le nom, déjà ?}
\phr{Sorry --- the line is breaking up.}{Pardon --- la ligne est mauvaise.}
\phr{Could you repeat the number, please?}{Pouvez-vous répéter le numéro ?}
\phr{Thanks for your patience with my English.}{Merci de votre patience avec mon anglais.}
\end{longtable}
\end{tableau}

\section{5. Donner et demander une opinion}
\begin{tableau}[30 phrases]
\footnotesize
\begin{longtable}{@{}p{8.0cm}p{6.6cm}@{}}
\entete{Anglais} & \entete{Français} \tabularnewline
\endfirsthead
\entete{Anglais} & \entete{Français} \tabularnewline
\endhead
\phr{In my opinion, it's not a good idea.}{À mon avis, ce n'est pas une bonne idée.}
\phr{As far as I'm concerned, it works fine.}{En ce qui me concerne, ça marche très bien.}
\phr{I think we should start earlier.}{Je pense que nous devrions commencer plus tôt.}
\phr{I honestly believe that\ldots{}}{Je crois sincèrement que\ldots{}}
\phr{If you ask me, it's too expensive.}{Si vous voulez mon avis, c'est trop cher.}
\phr{To my mind, the second option is better.}{Selon moi, la deuxième option est meilleure.}
\phr{I'd say it depends on the situation.}{Je dirais que ça dépend de la situation.}
\phr{I'm not convinced, to be honest.}{Je ne suis pas convaincu, honnêtement.}
\phr{I have mixed feelings about it.}{J'ai des sentiments mitigés à ce sujet.}
\phr{I strongly disagree.}{Je suis fortement en désaccord.}
\phr{I couldn't agree more.}{Je ne pourrais pas être plus d'accord.}
\phr{That's exactly what I think.}{C'est exactement ce que je pense.}
\phr{You've got a point there.}{Vous marquez un point / vous avez raison sur ce point.}
\phr{I see it differently.}{Je vois les choses différemment.}
\phr{That's one way of looking at it.}{C'est une façon de voir les choses.}
\phr{What do you think about it?}{Qu'en pensez-vous ?}
\phr{What's your opinion on this?}{Quelle est votre opinion là-dessus ?}
\phr{How do you feel about the new rules?}{Que pensez-vous des nouvelles règles ?}
\phr{Do you agree with me?}{Êtes-vous d'accord avec moi ?}
\phr{Would you say it's worth it?}{Diriez-vous que ça vaut le coup ?}
\phr{Let me put it this way.}{Laissez-moi le formuler ainsi.}
\phr{The way I see it, we have two options.}{À mes yeux, nous avons deux options.}
\phr{It seems to me that nothing has changed.}{Il me semble que rien n'a changé.}
\phr{I might be wrong, but\ldots{}}{J'ai peut-être tort, mais\ldots{}}
\phr{Correct me if I'm wrong.}{Corrigez-moi si je me trompe.}
\phr{There's no doubt that education matters.}{Il ne fait aucun doute que l'éducation compte.}
\phr{It's obvious that prices are rising.}{Il est évident que les prix augmentent.}
\phr{Some people think that\ldots{}, but I disagree.}{Certains pensent que\ldots{}, mais je ne suis pas d'accord.}
\phr{On the one hand\ldots{}, on the other hand\ldots{}}{D'un côté\ldots{}, de l'autre\ldots{}}
\phr{All in all, I'd do it again.}{Tout bien considéré, je le referais.}
\end{longtable}
\end{tableau}

\section{6. Demander, offrir, proposer}
\begin{tableau}[30 phrases]
\footnotesize
\begin{longtable}{@{}p{8.0cm}p{6.6cm}@{}}
\entete{Anglais} & \entete{Français} \tabularnewline
\endfirsthead
\entete{Anglais} & \entete{Français} \tabularnewline
\endhead
\phr{Could you help me with this, please?}{Pourriez-vous m'aider avec ceci ?}
\phr{Would you mind opening the window?}{Cela vous ennuierait-il d'ouvrir la fenêtre ?}
\phr{Could you do me a favour?}{Pourriez-vous me rendre un service ?}
\phr{Do you have a minute?}{Avez-vous une minute ?}
\phr{I was wondering if you could\ldots{}}{Je me demandais si vous pouviez\ldots{}}
\phr{Would it be possible to change the date?}{Serait-il possible de changer la date ?}
\phr{Can I ask you something?}{Puis-je vous demander quelque chose ?}
\phr{Where can I find a pharmacy?}{Où puis-je trouver une pharmacie ?}
\phr{How do I get to the train station?}{Comment aller à la gare ?}
\phr{Is it far from here?}{Est-ce loin d'ici ?}
\phr{Could you show me on the map?}{Pouvez-vous me montrer sur la carte ?}
\phr{What time does it open?}{À quelle heure ça ouvre ?}
\phr{Can I have the menu, please?}{Puis-je avoir la carte, s'il vous plaît ?}
\phr{May I use your phone?}{Puis-je utiliser votre téléphone ?}
\phr{Can I pay by card?}{Puis-je payer par carte ?}
\phr{Could I have the bill, please?}{Pourrais-je avoir l'addition ?}
\phr{Let me help you with your bags.}{Laissez-moi vous aider avec vos sacs.}
\phr{Would you like a cup of tea?}{Voulez-vous une tasse de thé ?}
\phr{Can I get you anything?}{Puis-je vous apporter quelque chose ?}
\phr{Shall I close the door?}{Voulez-vous que je ferme la porte ?}
\phr{Do you want me to drive?}{Voulez-vous que je conduise ?}
\phr{I'll take care of it.}{Je m'en occupe.}
\phr{Why don't we meet on Friday?}{Et si on se voyait vendredi ?}
\phr{How about going for a walk?}{Que diriez-vous d'aller se promener ?}
\phr{Would you like to join us?}{Voulez-vous vous joindre à nous ?}
\phr{I'd like to invite you to dinner.}{J'aimerais vous inviter à dîner.}
\phr{Are you free this evening?}{Êtes-vous libre ce soir ?}
\phr{That's very kind of you, thanks.}{C'est très aimable à vous, merci.}
\phr{Don't mention it.}{Je vous en prie / il n'y a pas de quoi.}
\phr{Any time!}{Quand vous voulez !}
\end{longtable}
\end{tableau}

\section{7. Accepter, refuser, hésiter}
\begin{tableau}[28 phrases --- refuser poliment est une compétence]
\footnotesize
\begin{longtable}{@{}p{8.0cm}p{6.6cm}@{}}
\entete{Anglais} & \entete{Français} \tabularnewline
\endfirsthead
\entete{Anglais} & \entete{Français} \tabularnewline
\endhead
\phr{Yes, of course.}{Oui, bien sûr.}
\phr{That sounds perfect.}{Ça me paraît parfait.}
\phr{I'd love to, thanks!}{Avec plaisir, merci !}
\phr{Sure, why not.}{Bien sûr, pourquoi pas.}
\phr{Count me in.}{Comptez sur moi / je suis partant.}
\phr{That works for me.}{Cela me convient.}
\phr{Let's do it.}{Faisons-le.}
\phr{I'm in.}{J'en suis.}
\phr{I'm afraid I can't.}{J'ai bien peur de ne pas pouvoir.}
\phr{I'd love to, but I already have plans.}{J'adorerais, mais j'ai déjà quelque chose de prévu.}
\phr{Sorry, that won't be possible.}{Désolé, ce ne sera pas possible.}
\phr{I'd rather not, if you don't mind.}{Je préférerais ne pas, si cela ne vous dérange pas.}
\phr{Maybe another time.}{Une autre fois, peut-être.}
\phr{I'm not sure --- let me think about it.}{Je ne suis pas sûr --- laissez-moi y réfléchir.}
\phr{I need to check my schedule first.}{Je dois d'abord vérifier mon emploi du temps.}
\phr{Could we postpone it to next week?}{Pourrions-nous reporter à la semaine prochaine ?}
\phr{I'll let you know as soon as I can.}{Je vous préviens dès que possible.}
\phr{It depends --- what time exactly?}{Ça dépend --- à quelle heure exactement ?}
\phr{To be honest, I'm not very keen on it.}{Honnêtement, ça ne m'enthousiasme pas.}
\phr{I'm not really into that kind of thing.}{Ce n'est pas vraiment mon truc.}
\phr{Thanks for asking, but no.}{Merci de proposer, mais non.}
\phr{I appreciate the offer, but I'll pass.}{J'apprécie l'offre, mais je passe mon tour.}
\phr{No, thank you --- I'm fine.}{Non, merci --- ça va.}
\phr{Sorry, I can't make it.}{Désolé, je ne peux pas venir.}
\phr{Something has come up.}{J'ai un imprévu.}
\phr{I have to cancel --- I'm really sorry.}{Je dois annuler --- je suis vraiment désolé.}
\phr{Let's say yes for now and confirm later.}{Disons oui pour l'instant et confirmons plus tard.}
\phr{If it's not too much trouble\ldots{}}{Si ce n'est pas trop de peine\ldots{}}
\end{longtable}
\end{tableau}

\section{8. Remercier et s'excuser}
\begin{tableau}[24 phrases]
\footnotesize
\begin{longtable}{@{}p{8.0cm}p{6.6cm}@{}}
\entete{Anglais} & \entete{Français} \tabularnewline
\endfirsthead
\entete{Anglais} & \entete{Français} \tabularnewline
\endhead
\phr{Thank you very much.}{Merci beaucoup.}
\phr{Thanks a lot for your help.}{Merci beaucoup pour votre aide.}
\phr{I really appreciate it.}{J'apprécie vraiment.}
\phr{That's very kind of you.}{C'est très gentil de votre part.}
\phr{I can't thank you enough.}{Je ne vous remercierai jamais assez.}
\phr{You've saved my life.}{Vous m'avez sauvé la vie.}
\phr{You're welcome.}{De rien / je vous en prie.}
\phr{My pleasure.}{Avec plaisir.}
\phr{Not at all.}{Pas du tout.}
\phr{No problem.}{Pas de problème.}
\phr{It was nothing.}{Ce n'était rien.}
\phr{I'm sorry --- it's my fault.}{Je suis désolé --- c'est ma faute.}
\phr{I apologise for the delay.}{Je vous présente mes excuses pour le retard.}
\phr{Sorry to bother you.}{Désolé de vous déranger.}
\phr{I shouldn't have said that.}{Je n'aurais pas dû dire cela.}
\phr{Please forgive me.}{Pardonnez-moi, je vous prie.}
\phr{It won't happen again.}{Cela ne se reproduira pas.}
\phr{I made a mistake and I take responsibility.}{J'ai fait une erreur et j'en assume la responsabilité.}
\phr{No need to apologise.}{Pas besoin de vous excuser.}
\phr{Forget about it --- it's fine.}{Oubliez --- ce n'est rien.}
\phr{I hope I didn't offend you.}{J'espère ne pas vous avoir offensé.}
\phr{That was rude of me --- sorry.}{C'était impoli de ma part --- pardon.}
\phr{Thanks for understanding.}{Merci de votre compréhension.}
\phr{I'm grateful for your support.}{Je suis reconnaissant de votre soutien.}
\end{longtable}
\end{tableau}

\section{9. Goûts, sentiments, émotions}
\begin{tableau}[26 phrases]
\footnotesize
\begin{longtable}{@{}p{8.0cm}p{6.6cm}@{}}
\entete{Anglais} & \entete{Français} \tabularnewline
\endfirsthead
\entete{Anglais} & \entete{Français} \tabularnewline
\endhead
\phr{I love this song.}{J'adore cette chanson.}
\phr{I'm really into African literature.}{Je suis passionné de littérature africaine.}
\phr{I'm crazy about football.}{Je suis fou de football.}
\phr{I quite like cooking.}{J'aime assez cuisiner.}
\phr{I don't really like horror films.}{Je n'aime pas trop les films d'horreur.}
\phr{I can't stand the noise.}{Je ne supporte pas le bruit.}
\phr{I hate being late.}{Je déteste être en retard.}
\phr{It's not really my thing.}{Ce n'est pas vraiment mon truc.}
\phr{I've never been interested in cars.}{Les voitures ne m'ont jamais intéressé.}
\phr{It grows on you.}{On s'y attache avec le temps.}
\phr{I'm so excited about this trip!}{Je suis tellement impatient pour ce voyage !}
\phr{I'm really looking forward to it.}{J'attends cela avec impatience.}
\phr{I'm a bit nervous, actually.}{Je suis un peu nerveux, en fait.}
\phr{I'm worried about the exam.}{Je m'inquiète pour l'examen.}
\phr{I'm afraid of flying.}{J'ai peur de l'avion.}
\phr{I feel so much better today.}{Je me sens beaucoup mieux aujourd'hui.}
\phr{I'm tired of waiting.}{J'en ai assez d'attendre.}
\phr{I'm fed up with this weather.}{J'en ai marre de ce temps.}
\phr{I'm proud of you.}{Je suis fier de toi.}
\phr{It made me laugh so much.}{Ça m'a beaucoup fait rire.}
\phr{I was really moved by the film.}{J'ai été très ému par le film.}
\phr{It's a bit disappointing.}{C'est un peu décevant.}
\phr{I'm relieved --- everything went well.}{Je suis soulagé --- tout s'est bien passé.}
\phr{I don't mind waiting.}{Cela ne me dérange pas d'attendre.}
\phr{It doesn't bother me at all.}{Cela ne me gêne pas du tout.}
\phr{I'm fine with that.}{Ça me va.}
\end{longtable}
\end{tableau}

\section{10. Le temps, les dates, les rendez-vous}
\begin{tableau}[28 phrases]
\footnotesize
\begin{longtable}{@{}p{8.0cm}p{6.6cm}@{}}
\entete{Anglais} & \entete{Français} \tabularnewline
\endfirsthead
\entete{Anglais} & \entete{Français} \tabularnewline
\endhead
\phr{What time is it?}{Quelle heure est-il ?}
\phr{It's half past nine.}{Il est neuf heures et demie.}
\phr{It's a quarter to eight.}{Il est huit heures moins le quart.}
\phr{What day is it today?}{Quel jour sommes-nous ?}
\phr{Today is Wednesday, the twelfth of March.}{Nous sommes mercredi douze mars.}
\phr{What's the date today?}{Quelle est la date d'aujourd'hui ?}
\phr{I was born on the fifth of May, 1995.}{Je suis né le 5 mai 1995.}
\phr{The meeting starts at ten sharp.}{La réunion commence à dix heures précises.}
\phr{Let's meet at eight in the evening.}{Rendez-vous à vingt heures.}
\phr{Can we move it to Friday?}{Pouvons-nous le déplacer à vendredi ?}
\phr{I'm busy on Monday afternoon.}{Je suis occupé lundi après-midi.}
\phr{Are you free next week?}{Êtes-vous libre la semaine prochaine ?}
\phr{I'll be there in ten minutes.}{J'arrive dans dix minutes.}
\phr{It takes about half an hour by bus.}{Cela prend environ une demi-heure en bus.}
\phr{I've been waiting for twenty minutes.}{J'attends depuis vingt minutes.}
\phr{We're running late.}{Nous sommes en retard.}
\phr{The deadline is next Monday.}{La date limite est lundi prochain.}
\phr{I do this twice a week.}{Je fais cela deux fois par semaine.}
\phr{Every other day, I go running.}{Un jour sur deux, je vais courir.}
\phr{I work from Monday to Friday.}{Je travaille du lundi au vendredi.}
\phr{It happens once in a while.}{Cela arrive de temps en temps.}
\phr{I haven't seen him for ages.}{Je ne l'ai pas vu depuis des lustres.}
\phr{A long time ago, I lived abroad.}{Il y a longtemps, j'habitais à l'étranger.}
\phr{These days, I study every evening.}{En ce moment, j'étudie tous les soirs.}
\phr{Soon, I hope.}{Bientôt, j'espère.}
\phr{Not yet --- maybe later.}{Pas encore --- peut-être plus tard.}
\phr{It's time to go.}{Il est temps de partir.}
\phr{The shop is open until nine.}{Le magasin est ouvert jusqu'à neuf heures.}
\end{longtable}
\end{tableau}

\section{11. Au téléphone et en visioconférence}
\begin{tableau}[26 phrases]
\footnotesize
\begin{longtable}{@{}p{8.0cm}p{6.6cm}@{}}
\entete{Anglais} & \entete{Français} \tabularnewline
\endfirsthead
\entete{Anglais} & \entete{Français} \tabularnewline
\endhead
\phr{Hello, this is Moussa speaking.}{Bonjour, Moussa à l'appareil.}
\phr{May I speak to Mrs Diop, please?}{Puis-je parler à Mme Diop, s'il vous plaît ?}
\phr{Who's calling, please?}{C'est de la part de qui ?}
\phr{Could you hold on a second?}{Pouvez-vous patienter un instant ?}
\phr{I'm afraid she's not available right now.}{Je crains qu'elle ne soit pas disponible actuellement.}
\phr{Can I take a message?}{Puis-je prendre un message ?}
\phr{Could you ask her to call me back?}{Pourriez-vous lui demander de me rappeler ?}
\phr{I'm calling about my order.}{J'appelle au sujet de ma commande.}
\phr{I'm returning your call.}{Je réponds à votre appel.}
\phr{Sorry, you've got the wrong number.}{Désolé, vous avez fait erreur de numéro.}
\phr{The line is bad --- can you hear me?}{La ligne est mauvaise --- vous m'entendez ?}
\phr{You're breaking up.}{Votre voix se coupe.}
\phr{Can you hear me now?}{Vous m'entendez maintenant ?}
\phr{I think we've been cut off.}{Je crois que nous avons été coupés.}
\phr{Let me put you through.}{Je vous passe quelqu'un.}
\phr{I'll call you back in five minutes.}{Je vous rappelle dans cinq minutes.}
\phr{Can everyone hear me all right?}{Tout le monde m'entend bien ?}
\phr{You're on mute.}{Vous êtes en mode muet.}
\phr{Could you share your screen, please?}{Pourriez-vous partager votre écran ?}
\phr{Sorry, go ahead --- you were speaking.}{Pardon, allez-y --- vous parliez.}
\phr{Let's schedule a call for Thursday.}{Programmons un appel pour jeudi.}
\phr{I'll send you the link by email.}{Je vous envoie le lien par courriel.}
\phr{My connection is unstable today.}{Ma connexion est instable aujourd'hui.}
\phr{Thanks for calling. Goodbye.}{Merci d'avoir appelé. Au revoir.}
\phr{I'll text you the details.}{Je vous envoie les détails par message.}
\phr{Speak to you soon!}{À bientôt au téléphone !}
\end{longtable}
\end{tableau}

\section{12. Au restaurant}
\begin{tableau}[26 phrases]
\footnotesize
\begin{longtable}{@{}p{8.0cm}p{6.6cm}@{}}
\entete{Anglais} & \entete{Français} \tabularnewline
\endfirsthead
\entete{Anglais} & \entete{Français} \tabularnewline
\endhead
\phr{A table for two, please.}{Une table pour deux, s'il vous plaît.}
\phr{Do you have a table by the window?}{Avez-vous une table près de la fenêtre ?}
\phr{What do you recommend?}{Que recommandez-vous ?}
\phr{What's in this dish?}{Qu'y a-t-il dans ce plat ?}
\phr{Is it spicy?}{Est-ce épicé ?}
\phr{I'm allergic to peanuts.}{Je suis allergique aux arachides.}
\phr{I don't eat meat --- I'm vegetarian.}{Je ne mange pas de viande --- je suis végétarien.}
\phr{Could I have this without onions?}{Pourrais-je l'avoir sans oignons ?}
\phr{I'll have the grilled fish, please.}{Je prendrai le poisson grillé, s'il vous plaît.}
\phr{Same for me.}{Pareil pour moi.}
\phr{Could we have some water, please?}{Pourrions-nous avoir de l'eau ?}
\phr{Still or sparkling?}{Plate ou gazeuse ?}
\phr{A glass of red wine, please.}{Un verre de vin rouge, s'il vous plaît.}
\phr{This is delicious!}{C'est délicieux !}
\phr{It's a bit cold, I'm afraid.}{C'est un peu froid, j'ai peur.}
\phr{I think there's been a mistake --- this isn't what I ordered.}{Je crois qu'il y a une erreur --- ce n'est pas ce que j'ai commandé.}
\phr{Could we have some more bread?}{Pourrions-nous avoir encore du pain ?}
\phr{Is service included?}{Le service est-il compris ?}
\phr{Could we have the bill, please?}{L'addition, s'il vous plaît.}
\phr{Can we pay separately?}{Pouvons-nous payer séparément ?}
\phr{Keep the change.}{Gardez la monnaie.}
\phr{The meal was excellent --- thank you.}{Le repas était excellent --- merci.}
\phr{Could I get this to take away?}{Puis-je l'emporter ?}
\phr{We're in a hurry --- is it quick to prepare?}{Nous sommes pressés --- est-ce rapide à préparer ?}
\phr{Do you accept credit cards?}{Acceptez-vous les cartes bancaires ?}
\phr{That was lovely --- we'll come back.}{C'était très agréable --- nous reviendrons.}
\end{longtable}
\end{tableau}

\section{13. À l'hôtel, en voyage, dans les transports}
\begin{tableau}[30 phrases]
\footnotesize
\begin{longtable}{@{}p{8.0cm}p{6.6cm}@{}}
\entete{Anglais} & \entete{Français} \tabularnewline
\endfirsthead
\entete{Anglais} & \entete{Français} \tabularnewline
\endhead
\phr{I have a reservation under the name Diop.}{J'ai une réservation au nom de Diop.}
\phr{I'd like to book a double room for two nights.}{Je voudrais réserver une chambre double pour deux nuits.}
\phr{Is breakfast included?}{Le petit-déjeuner est-il inclus ?}
\phr{What time is check-out?}{À quelle heure faut-il libérer la chambre ?}
\phr{Could I have a wake-up call at six?}{Pourrais-je avoir un réveil téléphonique à six heures ?}
\phr{The air conditioning isn't working.}{La climatisation ne fonctionne pas.}
\phr{Could I change rooms, please?}{Pourrais-je changer de chambre ?}
\phr{Is there Wi-Fi in the hotel?}{Y a-t-il du Wi-Fi dans l'hôtel ?}
\phr{What's the password, please?}{Quel est le mot de passe ?}
\phr{Could you call a taxi for me?}{Pourriez-vous m'appeler un taxi ?}
\phr{Where is the nearest bus stop?}{Où est l'arrêt de bus le plus proche ?}
\phr{Which platform does the train leave from?}{De quel quai part le train ?}
\phr{Is this seat taken?}{Cette place est-elle libre ?}
\phr{One ticket to London, please.}{Un billet pour Londres, s'il vous plaît.}
\phr{A return ticket --- how much is it?}{Un aller-retour --- combien cela coûte-t-il ?}
\phr{Is there a discount for students?}{Y a-t-il une réduction pour les étudiants ?}
\phr{What time does the flight leave?}{À quelle heure décolle le vol ?}
\phr{I'd like to check in, please.}{Je voudrais m'enregistrer, s'il vous plaît.}
\phr{Here's my passport and my ticket.}{Voici mon passeport et mon billet.}
\phr{I have one suitcase to check in.}{J'ai une valise à enregistrer.}
\phr{Is this hand luggage allowed?}{Ce bagage à main est-il autorisé ?}
\phr{My flight has been delayed.}{Mon vol a été retardé.}
\phr{I've missed my connection.}{J'ai raté ma correspondance.}
\phr{Where can I find the baggage claim?}{Où se trouve la livraison des bagages ?}
\phr{My suitcase hasn't arrived.}{Ma valise n'est pas arrivée.}
\phr{Could you help me, please? I'm lost.}{Pouvez-vous m'aider ? Je suis perdu.}
\phr{How long does it take to get there?}{Combien de temps faut-il pour y arriver ?}
\phr{Turn left at the second street.}{Tournez à gauche à la deuxième rue.}
\phr{Is it within walking distance?}{Est-ce accessible à pied ?}
\phr{Thanks --- you've been very helpful.}{Merci --- vous m'avez été très utile.}
\end{longtable}
\end{tableau}

\section{14. Achats et argent}
\begin{tableau}[24 phrases]
\footnotesize
\begin{longtable}{@{}p{8.0cm}p{6.6cm}@{}}
\entete{Anglais} & \entete{Français} \tabularnewline
\endfirsthead
\entete{Anglais} & \entete{Français} \tabularnewline
\endhead
\phr{I'm just looking, thanks.}{Je regarde seulement, merci.}
\phr{Do you have this in a larger size?}{Avez-vous ceci en taille plus grande ?}
\phr{Can I try it on?}{Puis-je l'essayer ?}
\phr{Where are the fitting rooms?}{Où sont les cabines d'essayage ?}
\phr{It doesn't fit --- it's too tight.}{Ça ne va pas --- c'est trop serré.}
\phr{How much is this?}{Combien coûte ceci ?}
\phr{That's a bit expensive for me.}{C'est un peu cher pour moi.}
\phr{Is there any discount?}{Y a-t-il une réduction ?}
\phr{Can you give me a better price?}{Pouvez-vous me faire un meilleur prix ?}
\phr{It's a good deal.}{C'est une bonne affaire.}
\phr{I'll take it.}{Je le prends.}
\phr{Could I have a receipt, please?}{Pourrais-je avoir un reçu ?}
\phr{Do you deliver?}{Livrez-vous ?}
\phr{Is it in stock?}{Est-ce en stock ?}
\phr{When will it be available again?}{Quand sera-t-il de nouveau disponible ?}
\phr{I'd like to return this --- it's faulty.}{Je voudrais retourner ceci --- il est défectueux.}
\phr{Can I get a refund or an exchange?}{Puis-je obtenir un remboursement ou un échange ?}
\phr{Here's my receipt.}{Voici mon reçu.}
\phr{Cash or card?}{En espèces ou par carte ?}
\phr{Could you break this note for me?}{Pouvez-vous me faire de la monnaie sur ce billet ?}
\phr{It costs fifty dollars.}{Cela coûte cinquante dollars.}
\phr{That's ten pounds each.}{C'est dix livres pièce.}
\phr{It's half price today.}{C'est à moitié prix aujourd'hui.}
\phr{Thanks for your help.}{Merci pour votre aide.}
\end{longtable}
\end{tableau}

\section{15. Santé et urgences}
\begin{tableau}[22 phrases]
\footnotesize
\begin{longtable}{@{}p{8.0cm}p{6.6cm}@{}}
\entete{Anglais} & \entete{Français} \tabularnewline
\endfirsthead
\entete{Anglais} & \entete{Français} \tabularnewline
\endhead
\phr{I don't feel well.}{Je ne me sens pas bien.}
\phr{I have a headache and a fever.}{J'ai mal à la tête et de la fièvre.}
\phr{I've had a sore throat for three days.}{J'ai mal à la gorge depuis trois jours.}
\phr{I've been coughing since Monday.}{Je tousse depuis lundi.}
\phr{My stomach hurts.}{J'ai mal au ventre.}
\phr{I think I've caught a cold.}{Je crois que j'ai attrapé un rhume.}
\phr{I'm allergic to penicillin.}{Je suis allergique à la pénicilline.}
\phr{I'm taking medication for my heart.}{Je prends un traitement pour le c\oe ur.}
\phr{I need to see a doctor.}{J'ai besoin de voir un médecin.}
\phr{Where is the nearest hospital?}{Où est l'hôpital le plus proche ?}
\phr{I'd like to make an appointment.}{Je voudrais prendre rendez-vous.}
\phr{Is it serious?}{Est-ce grave ?}
\phr{How long have you had these symptoms?}{Depuis combien de temps avez-vous ces symptômes ?}
\phr{Take this twice a day after meals.}{Prenez ceci deux fois par jour après les repas.}
\phr{You should rest for a few days.}{Vous devriez vous reposer quelques jours.}
\phr{Call an ambulance!}{Appelez une ambulance !}
\phr{Help! Somebody call the police!}{Au secours ! Que quelqu'un appelle la police !}
\phr{There's been an accident.}{Il y a eu un accident.}
\phr{I've lost my passport.}{J'ai perdu mon passeport.}
\phr{My bag has been stolen.}{On m'a volé mon sac.}
\phr{Where is the police station?}{Où est le commissariat ?}
\phr{I need to contact my embassy.}{Je dois contacter mon ambassade.}
\end{longtable}
\end{tableau}

\section{16. Au travail et en réunion}
\begin{tableau}[30 phrases]
\footnotesize
\begin{longtable}{@{}p{8.0cm}p{6.6cm}@{}}
\entete{Anglais} & \entete{Français} \tabularnewline
\endfirsthead
\entete{Anglais} & \entete{Français} \tabularnewline
\endhead
\phr{Let's get started, shall we?}{Commençons, voulez-vous ?}
\phr{The purpose of this meeting is to\ldots{}}{L'objet de cette réunion est de\ldots{}}
\phr{I'd like to go over three points.}{Je voudrais passer en revue trois points.}
\phr{First of all, let's look at the figures.}{Tout d'abord, regardons les chiffres.}
\phr{Could you give us an update on the project?}{Pouvez-vous nous faire un point sur le projet ?}
\phr{What's the current status?}{Où en sommes-nous ?}
\phr{I'd like to add something.}{J'aimerais ajouter quelque chose.}
\phr{Sorry to interrupt, but\ldots{}}{Pardon d'interrompre, mais\ldots{}}
\phr{If I may come back to what you said\ldots{}}{Si je puis revenir sur ce que vous avez dit\ldots{}}
\phr{Could you clarify that point?}{Pourriez-vous clarifier ce point ?}
\phr{What do you mean by ``efficiency''?}{Qu'entendez-vous par « efficacité » ?}
\phr{That's a good question --- let me think.}{C'est une bonne question --- laissez-moi réfléchir.}
\phr{I'll get back to you on that.}{Je reviens vers vous à ce sujet.}
\phr{Let's not go off topic.}{Ne nous écartons pas du sujet.}
\phr{Can we come back to the main point?}{Pouvons-nous revenir au point principal ?}
\phr{I'm afraid I don't have the figures with me.}{Je crains de ne pas avoir les chiffres avec moi.}
\phr{We're running out of time.}{Nous manquons de temps.}
\phr{Let's leave that for the next meeting.}{Laissons cela pour la prochaine réunion.}
\phr{To sum up, we agree on two points.}{Pour résumer, nous sommes d'accord sur deux points.}
\phr{So the next step is to contact the client.}{La prochaine étape est donc de contacter le client.}
\phr{Who's going to be in charge of this?}{Qui va s'en charger ?}
\phr{I'll take care of the report.}{Je m'occupe du rapport.}
\phr{Could you send me the minutes?}{Pourriez-vous m'envoyer le compte rendu ?}
\phr{The deadline is the end of the month.}{L'échéance est la fin du mois.}
\phr{We're behind schedule.}{Nous sommes en retard sur le planning.}
\phr{We've made good progress.}{Nous avons bien progressé.}
\phr{I'll keep you informed.}{Je vous tiendrai informé.}
\phr{Thanks everyone --- see you next week.}{Merci à tous --- à la semaine prochaine.}
\phr{I'm applying for the position of manager.}{Je postule au poste de responsable.}
\phr{Could I take a day off on Friday?}{Puis-je prendre un jour de congé vendredi ?}
\end{longtable}
\end{tableau}

\section{17. Apprendre une langue et parler de soi en tant qu'apprenant}
\begin{tableau}[24 phrases]
\footnotesize
\begin{longtable}{@{}p{8.0cm}p{6.6cm}@{}}
\entete{Anglais} & \entete{Français} \tabularnewline
\endfirsthead
\entete{Anglais} & \entete{Français} \tabularnewline
\endhead
\phr{I've been learning English for one year.}{J'apprends l'anglais depuis un an.}
\phr{My English is getting better every day.}{Mon anglais s'améliore chaque jour.}
\phr{I still translate in my head sometimes.}{Il m'arrive encore de traduire dans ma tête.}
\phr{Speaking is harder for me than writing.}{Parler est plus difficile pour moi qu'écrire.}
\phr{I understand much more than I can say.}{Je comprends bien plus que je ne peux dire.}
\phr{Could you correct my pronunciation?}{Pouvez-vous corriger ma prononciation ?}
\phr{Please correct me when I make mistakes.}{Corrigez-moi quand je fais des erreurs.}
\phr{Is this sentence natural?}{Cette phrase est-elle naturelle ?}
\phr{Would a native speaker say it like that?}{Un natif le dirait-il ainsi ?}
\phr{What's the difference between ``say'' and ``tell''?}{Quelle est la différence entre « say » et « tell » ?}
\phr{How do you pronounce this word?}{Comment prononce-t-on ce mot ?}
\phr{Is ``borrow'' the same as ``lend''?}{« Borrow » est-il identique à « lend » ?}
\phr{I need more practice with listening.}{J'ai besoin de plus d'entraînement à l'écoute.}
\phr{I watch one episode a day in English.}{Je regarde un épisode par jour en anglais.}
\phr{I keep a vocabulary notebook.}{Je tiens un carnet de vocabulaire.}
\phr{I review my flashcards every morning.}{Je révise mes cartes mémoire chaque matin.}
\phr{I'd like to find a conversation partner.}{J'aimerais trouver un partenaire de conversation.}
\phr{Do you know anyone who wants to practise?}{Connaissez-vous quelqu'un qui veut s'entraîner ?}
\phr{Could we speak English for half an hour?}{Pourrions-nous parler anglais pendant une demi-heure ?}
\phr{Let's switch to English, if that's all right.}{Passons à l'anglais, si cela vous convient.}
\phr{I'm preparing for the TOEIC exam.}{Je prépare l'examen TOEIC.}
\phr{My goal is to work in an international company.}{Mon objectif est de travailler dans une entreprise internationale.}
\phr{I'm proud of the progress I've made.}{Je suis fier des progrès que j'ai faits.}
\phr{One day, I'll think in English.}{Un jour, je penserai en anglais.}
\end{longtable}
\end{tableau}

\section{18. Conversation sociale (\emph{small talk})}
\begin{tableau}[28 phrases --- l'art de remplir les silences]
\footnotesize
\begin{longtable}{@{}p{8.0cm}p{6.6cm}@{}}
\entete{Anglais} & \entete{Français} \tabularnewline
\endfirsthead
\entete{Anglais} & \entete{Français} \tabularnewline
\endhead
\phr{Lovely weather today, isn't it?}{Il fait beau aujourd'hui, n'est-ce pas ?}
\phr{It's been raining for a week.}{Il pleut depuis une semaine.}
\phr{Is it always this hot in July?}{Fait-il toujours aussi chaud en juillet ?}
\phr{Do you live around here?}{Vous habitez dans le coin ?}
\phr{How long have you been in this city?}{Depuis combien de temps êtes-vous dans cette ville ?}
\phr{What do you do for a living?}{Que faites-vous dans la vie ?}
\phr{How did you get into that field?}{Comment en êtes-vous arrivé à ce métier ?}
\phr{Do you have any brothers or sisters?}{Avez-vous des frères ou s\oe urs ?}
\phr{What do you do in your free time?}{Que faites-vous pendant votre temps libre ?}
\phr{Have you seen any good films recently?}{Avez-vous vu de bons films récemment ?}
\phr{I'd recommend it --- it's excellent.}{Je vous le recommande --- c'est excellent.}
\phr{What kind of music do you like?}{Quel genre de musique aimez-vous ?}
\phr{Do you play any sports?}{Faites-vous du sport ?}
\phr{Did you support any team at the World Cup?}{Avez-vous soutenu une équipe à la Coupe du monde ?}
\phr{Have you ever been abroad?}{Êtes-vous déjà allé à l'étranger ?}
\phr{Where would you love to travel?}{Où aimeriez-vous voyager ?}
\phr{I've never been to Europe, but I'd love to go.}{Je ne suis jamais allé en Europe, mais j'adorerais y aller.}
\phr{What do you think of the food here?}{Que pensez-vous de la nourriture ici ?}
\phr{Have you tried the local dishes?}{Avez-vous goûté les plats locaux ?}
\phr{How was your weekend?}{Comment s'est passé votre week-end ?}
\phr{Nothing special --- I stayed at home.}{Rien de particulier --- je suis resté chez moi.}
\phr{What are your plans for the holidays?}{Quels sont vos projets pour les vacances ?}
\phr{That sounds like fun.}{Ça a l'air amusant.}
\phr{Anyway, it was nice chatting with you.}{En tout cas, c'était agréable de bavarder avec vous.}
\phr{Let's exchange numbers.}{Échangeons nos numéros.}
\phr{Are you on WhatsApp?}{Êtes-vous sur WhatsApp ?}
\phr{We should do this again.}{Il faudrait recommencer.}
\phr{Take care of yourself.}{Prenez soin de vous.}
\end{longtable}
\end{tableau}

\section{19. Raconter une histoire}
\begin{tableau}[24 phrases et connecteurs de récit]
\footnotesize
\begin{longtable}{@{}p{8.0cm}p{6.6cm}@{}}
\entete{Anglais} & \entete{Français} \tabularnewline
\endfirsthead
\entete{Anglais} & \entete{Français} \tabularnewline
\endhead
\phr{Let me tell you what happened.}{Laissez-moi vous raconter ce qui s'est passé.}
\phr{It all started last summer.}{Tout a commencé l'été dernier.}
\phr{One day, I was walking home when\ldots{}}{Un jour, je rentrais chez moi quand\ldots{}}
\phr{At first, I didn't understand.}{Au début, je n'ai pas compris.}
\phr{Then something strange happened.}{Puis quelque chose d'étrange s'est produit.}
\phr{Suddenly, the lights went out.}{Soudain, les lumières se sont éteintes.}
\phr{A few minutes later, he arrived.}{Quelques minutes plus tard, il est arrivé.}
\phr{Meanwhile, my sister was waiting outside.}{Pendant ce temps, ma s\oe ur attendait dehors.}
\phr{In the end, everything was fine.}{Finalement, tout s'est bien passé.}
\phr{It turned out that he knew her.}{Il s'est avéré qu'il la connaissait.}
\phr{To cut a long story short\ldots{}}{Pour faire court\ldots{}}
\phr{The funny thing is that\ldots{}}{Le plus drôle, c'est que\ldots{}}
\phr{Believe it or not, I won.}{Croyez-le ou non, j'ai gagné.}
\phr{I'll never forget that day.}{Je n'oublierai jamais ce jour-là.}
\phr{Looking back, it was a great experience.}{Avec le recul, ce fut une belle expérience.}
\phr{That's when I realised something important.}{C'est là que j'ai compris quelque chose d'important.}
\phr{If it hadn't been for him, I would have failed.}{Sans lui, j'aurais échoué.}
\phr{I used to do this every Sunday.}{Je faisais cela tous les dimanches.}
\phr{We were about to leave when it started.}{Nous étions sur le point de partir quand cela a commencé.}
\phr{He kept talking for an hour.}{Il n'a pas arrêté de parler pendant une heure.}
\phr{Nobody expected that.}{Personne ne s'y attendait.}
\phr{And that's the story.}{Et voilà l'histoire.}
\phr{What happened next?}{Que s'est-il passé ensuite ?}
\phr{How did you feel at that moment?}{Comment vous êtes-vous senti à ce moment-là ?}
\end{longtable}
\end{tableau}

\section{20. Parler de ses projets et de l'avenir}
\begin{tableau}[22 phrases]
\footnotesize
\begin{longtable}{@{}p{8.0cm}p{6.6cm}@{}}
\entete{Anglais} & \entete{Français} \tabularnewline
\endfirsthead
\entete{Anglais} & \entete{Français} \tabularnewline
\endhead
\phr{I'm going to start a new job in June.}{Je vais commencer un nouveau travail en juin.}
\phr{I'm planning to travel next year.}{Je prévois de voyager l'an prochain.}
\phr{I hope to improve my English quickly.}{J'espère améliorer mon anglais rapidement.}
\phr{My goal is to become fluent by December.}{Mon objectif est de devenir à l'aise d'ici décembre.}
\phr{I'd like to work for an international company.}{J'aimerais travailler pour une entreprise internationale.}
\phr{I'm thinking of going back to university.}{J'envisage de retourner à l'université.}
\phr{One day, I'd love to live abroad.}{Un jour, j'adorerais vivre à l'étranger.}
\phr{If everything goes well, I'll move in the spring.}{Si tout va bien, je déménagerai au printemps.}
\phr{I might change career --- nothing is decided.}{Je changerai peut-être de carrière --- rien n'est décidé.}
\phr{I've just started my own business.}{Je viens de créer ma propre entreprise.}
\phr{I'm saving money for a project.}{J'économise pour un projet.}
\phr{In five years, I see myself\ldots{}}{Dans cinq ans, je me vois\ldots{}}
\phr{I don't know yet, but I'll figure it out.}{Je ne sais pas encore, mais je trouverai.}
\phr{It depends on how things go.}{Cela dépend de comment les choses évoluent.}
\phr{I've made up my mind.}{J'ai pris ma décision.}
\phr{I'm still weighing up the options.}{Je pèse encore les options.}
\phr{Let's see how it goes.}{On verra comment ça se passe.}
\phr{I'm looking forward to the next step.}{J'attends la prochaine étape avec impatience.}
\phr{Whatever happens, I'll keep learning.}{Quoi qu'il arrive, je continuerai à apprendre.}
\phr{This is only the beginning.}{Ce n'est que le début.}
\phr{I believe hard work always pays off.}{Je crois que le travail finit toujours par payer.}
\phr{Thank you --- you've been very kind.}{Merci --- vous avez été très aimable.}
\end{longtable}
\end{tableau}

\section{21. Phrases pour l'écrit : e-mails et messages}
\begin{tableau}[32 formules d'e-mail --- du plus formel au plus courant]
\footnotesize
\begin{longtable}{@{}p{8.0cm}p{6.6cm}@{}}
\entete{Anglais} & \entete{Français} \tabularnewline
\endfirsthead
\entete{Anglais} & \entete{Français} \tabularnewline
\endhead
\phr{Dear Sir or Madam,}{Madame, Monsieur,}
\phr{Dear Mr Diop,}{Cher Monsieur Diop,}
\phr{Dear Aïcha,}{Chère Aïcha,}
\phr{Hi Sam --- hope you're doing well.}{Salut Sam --- j'espère que tu vas bien.}
\phr{I hope this email finds you well.}{J'espère que ce courriel vous trouvera en bonne santé.}
\phr{I'm writing to enquire about\ldots{}}{Je vous écris pour me renseigner au sujet de\ldots{}}
\phr{I'm writing in connection with your advertisement.}{Je vous écris au sujet de votre annonce.}
\phr{I'm writing to follow up on our conversation.}{Je vous écris pour faire suite à notre échange.}
\phr{Further to our phone call this morning\ldots{}}{Suite à notre appel téléphonique de ce matin\ldots{}}
\phr{With reference to your email of 3 May\ldots{}}{En référence à votre courriel du 3 mai\ldots{}}
\phr{I would like to request more information.}{Je souhaiterais obtenir davantage d'informations.}
\phr{Could you please send me the document?}{Pourriez-vous m'envoyer le document ?}
\phr{I would be grateful if you could reply by Friday.}{Je vous serais reconnaissant de bien vouloir répondre avant vendredi.}
\phr{Please find attached my CV.}{Veuillez trouver ci-joint mon CV.}
\phr{I've attached the report to this email.}{J'ai joint le rapport à ce courriel.}
\phr{Just a quick note to say thanks.}{Un petit mot pour vous remercier.}
\phr{Sorry for the late reply.}{Désolé pour la réponse tardive.}
\phr{Apologies for the delay in getting back to you.}{Mes excuses pour le délai de ma réponse.}
\phr{Thank you for your prompt reply.}{Merci pour votre réponse rapide.}
\phr{That's very clear --- thank you.}{C'est très clair --- merci.}
\phr{I'm afraid there's been a misunderstanding.}{Je crains qu'il y ait eu un malentendu.}
\phr{Could you clarify one point, please?}{Pourriez-vous clarifier un point, s'il vous plaît ?}
\phr{Unfortunately, we are unable to accept your request.}{Malheureusement, nous ne sommes pas en mesure d'accepter votre demande.}
\phr{We will do our best to solve the problem.}{Nous ferons de notre mieux pour résoudre le problème.}
\phr{Please don't hesitate to contact me.}{N'hésitez pas à me contacter.}
\phr{Let me know if you need anything else.}{Dites-moi si vous avez besoin d'autre chose.}
\phr{I look forward to hearing from you.}{Dans l'attente de votre réponse.}
\phr{Looking forward to meeting you.}{Au plaisir de vous rencontrer.}
\phr{Kind regards,}{Cordialement,}
\phr{Best regards,}{Bien cordialement,}
\phr{Yours sincerely,}{Veuillez agréer, Madame, Monsieur, mes salutations distinguées,}
\phr{See you soon! Talk later.}{À bientôt ! On se parle plus tard.}
\end{longtable}
\end{tableau}

\section{22. Tournures naturelles qui font « anglais courant »}
\begin{tableau}[30 expressions --- celles qui font la différence]
\footnotesize
\begin{longtable}{@{}p{8.0cm}p{6.6cm}@{}}
\entete{Anglais} & \entete{Français} \tabularnewline
\endfirsthead
\entete{Anglais} & \entete{Français} \tabularnewline
\endhead
\phr{It's up to you.}{C'est à toi de décider.}
\phr{That's not my cup of tea.}{Ce n'est pas mon truc.}
\phr{Let's play it by ear.}{Improvisons, on verra bien.}
\phr{It's not a big deal.}{Ce n'est pas grave.}
\phr{No biggie.}{Pas de souci. (familier)}
\phr{Long story short, I lost my job.}{Bref, j'ai perdu mon travail.}
\phr{Make yourself at home.}{Faites comme chez vous.}
\phr{It's on me --- my treat.}{C'est pour moi --- je vous invite.}
\phr{Let's split the bill.}{Partageons l'addition.}
\phr{I'm on my way.}{J'arrive / je suis en route.}
\phr{I'll be right there.}{J'arrive tout de suite.}
\phr{Give me a hand, will you?}{Donne-moi un coup de main, tu veux ?}
\phr{I'm swamped at work this week.}{Je suis débordé au travail cette semaine.}
\phr{Take it easy.}{Vas-y doucement / détends-toi.}
\phr{Don't worry about it.}{Ne t'inquiète pas pour ça.}
\phr{It doesn't make any sense.}{Cela n'a aucun sens.}
\phr{I've had enough of this.}{J'en ai assez de ça.}
\phr{Let's call it a day.}{Arrêtons-nous là pour aujourd'hui.}
\phr{It's worth a try.}{Cela vaut la peine d'essayer.}
\phr{There's no point in arguing.}{Il est inutile de se disputer.}
\phr{I can't afford it right now.}{Je ne peux pas me le permettre actuellement.}
\phr{It cost me an arm and a leg.}{Cela m'a coûté les yeux de la tête.}
\phr{That's a rip-off!}{C'est du vol !}
\phr{I'll get to it later.}{Je m'en occuperai plus tard.}
\phr{It slipped my mind.}{Cela m'est sorti de la tête.}
\phr{I'm about to leave.}{Je suis sur le point de partir.}
\phr{It turned out better than expected.}{Cela s'est avéré mieux que prévu.}
\phr{Everything went smoothly.}{Tout s'est déroulé sans accroc.}
\phr{I'm getting the hang of it.}{Je commence à piger.}
\phr{Practice makes perfect.}{C'est en forgeant qu'on devient forgeron.}
\end{longtable}
\end{tableau}

\begin{mission}[Annexe B --- le plan des 40 jours]
600 phrases, 15 par jour, 40 jours. Chaque matin : 15 phrases à voix haute (10 minutes).
Chaque soir : utilisez-en au moins trois dans la journée --- en conversation réelle, en
enregistrement, ou simplement en vous parlant. Le dimanche : test des 105 phrases de la
semaine, colonne française cachée. Au jour 40, vous possédez un répertoire qui couvre
90\,\% des situations de la vie quotidienne. C'est le plus grand raccourci vers la fluidité :
\textbf{parler en blocs tout faits pendant que votre grammaire, elle, continue de se
construire.}
\end{mission}
